% =========================================================
% Characterization of Discontinuity — Notes
% Chapter: continuity
% Volume III · Analysis
% =========================================================

\section{Characterization of Discontinuity}
\label{sec:discontinuity}

\begin{tcolorbox}[title={Toolkit --- Discontinuity}]
\begin{itemize}
  \item $\operatorname{DiscontinuousAt}(f,a)$ $\iff$ negation of
    $\operatorname{ContinuousAt}(f,a)$. Three equivalent forms:
    $\varepsilon$-$\delta$, sequential, neighbourhood.
  \item Quantifier reversal from continuity:
    $\forall\varepsilon\leadsto\exists\varepsilon_0$,\;
    $\exists\delta\leadsto\forall\delta$,\;
    $\forall x\leadsto\exists x$.
  \item Taxonomy: removable, jump (first kind), essential (second kind).
  \item First kind: both one-sided limits $f(a^-)$, $f(a^+)$ exist finitely.
  \item Second kind: at least one one-sided limit fails to exist finitely.
  \item Every discontinuity is of exactly one kind.
\end{itemize}
\end{tcolorbox}

% ---------------------------------------------------------
% Definition — Point of Discontinuity
% ---------------------------------------------------------

\begin{tcolorbox}[title={Definition --- Point of Discontinuity}]
\begin{definition}[Point of Discontinuity]
\label{def:point-of-discontinuity}
$\operatorname{DiscontinuousAt}(f,a) \colonequiv
\neg\,\operatorname{ContinuousAt}(f,a)$.
\end{definition}
\end{tcolorbox}

\begin{remark*}[Standard quantified statement]
\[
  \operatorname{DiscontinuousAt}(f,a)
  \;\iff\;
  \exists\varepsilon_0>0,\;\forall\delta>0,\;\exists x\in E:\;
  \operatorname{Within}(x,a,\delta)
  \;\wedge\;
  \neg\,\operatorname{OutputTolerance}(f(x),f(a),\varepsilon_0).
\]
\end{remark*}

\begin{remark*}[Definition predicate reading]
\[
  \operatorname{DiscontinuousAt}(f,a)
  \;\colonequiv\;
  \neg\,\operatorname{ContinuousAt}(f,a).
\]
\end{remark*}

\begin{remark*}[Failure mode decomposition — quantifier reversal]
\[
  \underbrace{\forall\varepsilon>0}_{\text{cts}}
  \;\leadsto\;
  \underbrace{\exists\varepsilon_0>0}_{\text{dis}},
  \qquad
  \underbrace{\exists\delta>0}_{\text{cts}}
  \;\leadsto\;
  \underbrace{\forall\delta>0}_{\text{dis}},
  \qquad
  \underbrace{\forall x\in A}_{\text{cts}}
  \;\leadsto\;
  \underbrace{\exists x\in E}_{\text{dis}}.
\]
The surviving gap $\varepsilon_0>0$ is the witness to discontinuity.
\end{remark*}

\begin{remark*}[Interpretation]
A fixed output gap cannot be closed no matter how tightly the input
window shrinks. Continuity fails because nearby values cannot be
forced into an arbitrarily small band around $f(a)$.
\end{remark*}

% ---------------------------------------------------------
% Definition — epsilon-delta Characterization
% ---------------------------------------------------------

\begin{definition}[$\varepsilon$-$\delta$ Characterization of Discontinuity]
\label{def:epsilon-delta-discontinuity-at-a-point}
$f$ is \textbf{discontinuous at $x_0$} ($\varepsilon$-$\delta$) if
\[
  \exists\varepsilon_0>0,\;\forall\delta>0,\;\exists x\in E:\;
  |x-x_0|<\delta \;\wedge\; |f(x)-f(x_0)|\geq\varepsilon_0.
\]
\end{definition}

\begin{remark*}[Standard quantified statement]
\[
  \operatorname{DiscontinuousAt}_{\varepsilon\delta}(f,x_0)
  \;\colonequiv\;
  \exists\varepsilon_0>0,\;\forall\delta>0,\;\exists x\in E:\;
  \operatorname{Within}(x,x_0,\delta)
  \;\wedge\;
  \neg\,\operatorname{OutputTolerance}(f(x),f(x_0),\varepsilon_0).
\]
\end{remark*}

\begin{remark*}[Interpretation]
Some fixed output gap survives every shrinking of the input window.
\end{remark*}

% ---------------------------------------------------------
% Definition — Sequential Characterization
% ---------------------------------------------------------

\begin{definition}[Sequential Characterization of Discontinuity]
\label{def:sequential-discontinuity-at-a-point}
$f$ is \textbf{discontinuous at $x_0$} (sequential) if
\[
  \exists(x_n)\subseteq E:\;
  x_n\to x_0 \;\wedge\; f(x_n)\not\to f(x_0).
\]
\end{definition}

\begin{remark*}[Standard quantified statement]
\[
  \operatorname{DiscontinuousAt}_{\mathrm{seq}}(f,x_0)
  \;\colonequiv\;
  \exists(x_n)\subseteq E:\;
  \operatorname{ConvergentSequence}(x_n,x_0)
  \;\wedge\;
  \neg\,\operatorname{ConvergentSequence}(f(x_n),f(x_0)).
\]
\end{remark*}

\begin{remark*}[Interpretation]
One witness sequence converging to $x_0$ whose image fails to converge
to $f(x_0)$ suffices. Continuity demands every such sequence succeed;
one failure is enough to establish discontinuity.
\end{remark*}

% ---------------------------------------------------------
% Definition — Neighbourhood Characterization
% ---------------------------------------------------------

\begin{definition}[Neighbourhood Characterization of Discontinuity]
\label{def:neighborhood-discontinuity-at-a-point}
$f$ is \textbf{discontinuous at $x_0$} (neighbourhood) if
\[
  \exists V\text{ nbhd of }f(x_0),\;\forall U\text{ nbhd of }x_0,\;
  \exists x\in U\cap E: f(x)\notin V.
\]
\end{definition}

\begin{remark*}[Standard quantified statement]
\[
  \operatorname{DiscontinuousAt}_{\mathrm{nbhd}}(f,x_0)
  \;\colonequiv\;
  \exists\varepsilon_0>0,\;\forall\delta>0,\;\exists x\in E:\;
  \operatorname{Within}(x,x_0,\delta)
  \;\wedge\;
  \neg\,\operatorname{Within}(f(x),f(x_0),\varepsilon_0).
\]
\end{remark*}

\begin{remark*}[Interpretation]
A fixed target neighbourhood around $f(x_0)$ that the image of every
domain neighbourhood fails to stay inside. Continuity would require
some domain neighbourhood to map entirely inside any prescribed target;
discontinuity is the failure for at least one target.
\end{remark*}

% ---------------------------------------------------------
% Lemma — Equivalence of Three Formulations
% ---------------------------------------------------------

\begin{lemma}[Equivalent Formulations of Discontinuity at a Point]
\label{lem:equivalent-discontinuity-at-a-point}
The $\varepsilon$-$\delta$, sequential, and neighbourhood forms of
discontinuity at $x_0$ are equivalent.

\smallskip\noindent
\hyperref[prf:discontinuity-equivalence]{\textit{Go to proof.}}
\end{lemma}

\begin{remark*}[Interpretation]
Each is the negation of the corresponding formulation of continuity,
and they are equivalent by the same cyclic argument.
\end{remark*}

% ---------------------------------------------------------
% Definition — Types of Discontinuity
% ---------------------------------------------------------

\begin{tcolorbox}[title={Definition --- Types of Discontinuity at a Point}]
\begin{definition}[Types of Discontinuity at a Point]
\label{def:types-of-discontinuity-at-a-point}
Let $x_0\in E$ be a point of discontinuity of $f$.
\begin{enumerate}[label=(\roman*)]
  \item \textbf{Removable}: $\operatorname{FunctionLimit}(f,x_0,L)$
    exists but $L\neq f(x_0)$.
  \item \textbf{Jump}: $\operatorname{LeftLimit}(f,x_0,L_-)$ and
    $\operatorname{RightLimit}(f,x_0,L_+)$ both exist but $L_-\neq L_+$.
  \item \textbf{Essential}: not removable.
  \item \textbf{Infinite}: at least one one-sided limit is $\pm\infty$.
  \item \textbf{Oscillatory}: neither removable, jump, nor infinite.
\end{enumerate}
\end{definition}
\end{tcolorbox}

\begin{remark*}[Standard quantified statement]
\[
\begin{aligned}
  \text{Removable} &\;\colonequiv\;
    \exists L\in\mathbb{R}:\operatorname{FunctionLimit}(f,x_0,L)\wedge f(x_0)\neq L. \\
  \text{Jump} &\;\colonequiv\;
    \exists L_-,L_+\in\mathbb{R}:
    \operatorname{LeftLimit}(f,x_0,L_-)\wedge
    \operatorname{RightLimit}(f,x_0,L_+)\wedge L_-\neq L_+.
\end{aligned}
\]
\end{remark*}

\begin{remark*}[Failure mode decomposition]
\[
  \operatorname{DiscontinuousAt}(f,x_0)
  \;\begin{cases}
    \text{removable} & \text{limit exists, wrong value} \\
    \text{jump} & \text{both one-sided limits exist, unequal} \\
    \text{essential}
      \;\begin{cases}
        \text{infinite} & \text{some limit is }\pm\infty \\
        \text{oscillatory} & \text{no relevant limit exists}
      \end{cases}
  \end{cases}
\]
\end{remark*}

\begin{remark*}[Canonical examples]
(i) Removable: $f(x)=(x^2-1)/(x-1)$, $f(1)=0$; limit at $1$ is $2$.
(ii) Jump: $f(x)=\operatorname{sgn}(x)$ at $0$.
(iii) Infinite: $f(x)=1/x^2$ at $0$.
(iv) Oscillatory: $f(x)=\sin(1/x)$ at $0$.
\end{remark*}

\begin{remark*}[Interpretation]
Removable discontinuities have correct limiting behavior but wrong
point value — redefining $f$ at $x_0$ alone repairs them. Essential
discontinuities cannot be repaired at a single point.
\end{remark*}

% ---------------------------------------------------------
% Definition — First and Second Kind
% ---------------------------------------------------------

\begin{tcolorbox}[title={Definition --- Discontinuity of First Kind}]
\begin{definition}[Discontinuity of First Kind]
\label{def:discontinuity-of-first-kind}
$\operatorname{FirstKind}(f,a)$: both $\operatorname{LeftLimit}(f,a,L_-)$
and $\operatorname{RightLimit}(f,a,L_+)$ exist in $\mathbb{R}$.
\end{definition}
\end{tcolorbox}

\begin{remark*}[Standard quantified statement]
\[
  \operatorname{FirstKind}(f,a)
  \;\colonequiv\;
  \exists L_-,L_+\in\mathbb{R}:
  \operatorname{LeftLimit}(f,a,L_-)
  \;\wedge\;
  \operatorname{RightLimit}(f,a,L_+).
\]
\end{remark*}

\begin{remark*}[Interpretation]
First-kind includes both removable (one-sided limits exist and are
equal, but differ from $f(a)$) and jump (one-sided limits exist but
are unequal). The defining feature is existence of both finite
one-sided limits.
\end{remark*}

\begin{tcolorbox}[title={Definition --- Discontinuity of Second Kind}]
\begin{definition}[Discontinuity of Second Kind]
\label{def:discontinuity-of-second-kind}
$\operatorname{SecondKind}(f,a)$: $\neg\,\operatorname{FirstKind}(f,a)$.
At least one of $\operatorname{LeftLimit}(f,a,\cdot)$,
$\operatorname{RightLimit}(f,a,\cdot)$ does not exist as a real number.
\end{definition}
\end{tcolorbox}

\begin{remark*}[Standard quantified statement]
\[
  \operatorname{SecondKind}(f,a)
  \;\colonequiv\;
  \neg\,\operatorname{FirstKind}(f,a).
\]
\end{remark*>

\begin{remark*}[Interpretation]
Second-kind includes infinite blow-up and oscillatory behavior.
$\sin(1/x)$ at $0$ is the canonical second-kind example: oscillates
between $-1$ and $1$ in every neighbourhood of $0$, so no one-sided
limit exists. Every discontinuity is of exactly first or second kind.
\end{remark*}
